% =====================================================
% ROTEIRO PRÁTICO 03
% =====================================================

\section{Roteiro Prático 03 — Escrita Inicial do Artigo no Padrão ACM}

% =====================================================
% OBJETIVOS
% =====================================================

\begin{boxObjetivos}
\textbf{Objetivos da Aula}

Ao final desta aula o estudante será capaz de:

\begin{itemize}
\item Criar um artigo científico utilizando o template oficial ACM no Overleaf;
\item Configurar corretamente os dados iniciais do artigo;
\item Redigir o resumo e a introdução do projeto;
\item Iniciar a seção de fundamentação teórica (background);
\item Registrar o link do artigo na planilha de acompanhamento.
\end{itemize}
\end{boxObjetivos}

% =====================================================
% INTRODUÇÃO
% =====================================================

\subsection{Introdução}

A partir desta aula, os projetos passam a ser documentados no formato de artigo científico utilizando o padrão ACM, adotado em eventos como o Computer on the Beach.

O artigo será desenvolvido ao longo do semestre e poderá ser utilizado como base para submissão em evento científico.

O foco desta aula é iniciar a estrutura do artigo com os elementos essenciais.

% =====================================================
% PARTE 1
% =====================================================

\subsection{Parte 1 — Acesso ao Template ACM no Overleaf}

\begin{boxAtividade}[Atividade 1 — Criação do Projeto]

Cada grupo deverá:

\begin{enumerate}
\item Acessar o template oficial ACM no Overleaf:

\texttt{https://www.overleaf.com/gallery/tagged/acm-official\#.WOuOk2e1taQ}

\item Selecionar um template de conferência (conference);

\item Criar uma cópia do projeto no Overleaf;

\item Renomear o projeto com o nome do trabalho;

\item Verificar se o documento compila corretamente.
\end{enumerate}

\end{boxAtividade}

% =====================================================
% PARTE 2
% =====================================================

\subsection{Parte 2 — Configuração Inicial do Artigo}

\begin{boxAtividade}[Atividade 2 — Ajuste dos Dados]

No arquivo principal do artigo, atualizar:

\begin{itemize}
\item Título do trabalho;
\item Nome dos autores;
\item Instituição (IFSC – Campus Garopaba);
\item Remover textos de exemplo do template.
\end{itemize}

\end{boxAtividade}

% =====================================================
% PARTE 3
% =====================================================

\subsection{Parte 3 — Escrita do Resumo (Abstract)}

\begin{boxAtividade}[Atividade 3 — Resumo]

O resumo deve conter:

\begin{itemize}
\item O problema abordado;
\item A proposta do projeto;
\item O objetivo principal.
\end{itemize}

O texto deve ser breve, claro e direto (aproximadamente 5 a 8 linhas).

\end{boxAtividade}

% =====================================================
% PARTE 4
% =====================================================

\subsection{Parte 4 — Introdução}

\begin{boxAtividade}[Atividade 4 — Introdução]

A introdução deve responder:

\begin{itemize}
\item Qual problema o projeto resolve;
\item Por que esse problema é relevante;
\item Qual solução está sendo proposta.
\end{itemize}

A seção deve conter entre 2 e 4 parágrafos.

\end{boxAtividade}

% =====================================================
% PARTE 5
% =====================================================

\subsection{Parte 5 — Fundamentação Teórica (Background)}

\begin{boxAtividade}[Atividade 5 — Background/Fundamentação teórica]

Cada grupo deverá:

\begin{itemize}
\item Criar ou ajustar a seção de background;
\item Inserir pelo menos 2 referências;
\item Explicar conceitos relacionados ao projeto;
\item Relacionar esses conceitos com a solução proposta.
\end{itemize}

\end{boxAtividade}

% =====================================================
% PARTE 6
% =====================================================

\subsection{Parte 6 — Registro na Planilha}

\begin{boxAtividade}[Atividade 6 — Atualização do Acompanhamento]

Cada grupo deverá:

\begin{itemize}
\item Copiar o link do projeto no Overleaf;
\item Inserir o link na planilha de acompanhamento;
\item Atualizar o status do projeto;
\item Registrar a data da última atualização.
\end{itemize}

\end{boxAtividade}

% =====================================================
% CHECKLIST
% =====================================================

\subsection{Checklist Final da Aula}

\begin{boxAtividade}[Verificação Técnica]

Antes de finalizar:

\begin{itemize}
\item O artigo foi criado no Overleaf usando template ACM;
\item O título e autores foram atualizados;
\item O resumo foi escrito;
\item A introdução está estruturada;
\item A seção de background foi iniciada;
\item O link do artigo está na planilha.
\end{itemize}

\end{boxAtividade}
